%% ======================================================================
%% chap5/implementation_main.tex
%%
%% Chapter 5: Implementation
%%
%% Condensed implementation chapter. Development environment, toolchain,
%% workspace, and pallet internals (storage layout, extrinsic reference,
%% royalty mechanics, ink! wrapper, EVM precompile) are documented in
%% Appendix C. This chapter covers the pallet architecture at a summary
%% level, the XCM and Snowbridge cross-chain integration in full, and
%% the nine technical challenges that shaped the final design.
%%
%% Sections:
%%   5.1 Pallet Implementation and Smart Contract Layer (summary)
%%   5.2 Cross-Chain Integration (2 subsections)
%%   5.3 Challenges and Solutions (Table 5.1)
%%   5.4 Summary
%%
%% External labels referenced:
%%   chap:discussion (Ch 8), chap:conclusion (Ch 9),
%%   app:challenges (Appendix A), app:code (Appendix C),
%%   sec:contract-first-pivot (Appendix A.2),
%%   sec:children-bound (Appendix A.8),
%%   sec:ink-discontinuation (Appendix A.9)
%%
%% Citations: none (implementation narrative is self-contained).
%% ======================================================================

\chapter{Implementation}\label{chap:implementation}

This chapter describes what was built and how key decisions played out in practice. The development environment, toolchain, and workspace layout are documented in Appendix~\ref{app:code}. Section~\ref{sec:pallet-impl} summarizes the pallet implementation; Section~\ref{sec:cross-chain-integration} covers cross-chain integration; Section~\ref{sec:challenges-solutions} documents the nine challenges encountered and their resolutions.

\section{Pallet Implementation and Smart Contract Layer}\label{sec:pallet-impl}

Figure~\ref{fig:component-overview} shows the six implementation components and their relationships.

\begin{figure}[htbp]
\vspace{\baselineskip}
\centering
\begin{tikzpicture}[
    box/.style={draw, rectangle, rounded corners=4pt,
                minimum width=3.0cm, minimum height=1.0cm,
                align=center, font=\small},
    ext/.style={box, fill=gray!15},
    core/.style={box, fill=blue!20},
    dep/.style={box, fill=orange!15},
    sib/.style={box, fill=green!15},
    lbl/.style={font=\footnotesize, fill=white, inner sep=1.5pt},
    >=Stealth
]
%% Top row: external interface components
\node[ext] (xcm) at (-3.8, 3.8) {
    \textbf{XCM v3}\\
    {\footnotesize remote parachain calls}
};
\node[ext] (ink) at (0, 3.8) {
    \textbf{\texttt{rights\_manager}}\\
    {\footnotesize ink!\ 6 wrapper contract}
};
\node[ext] (pre) at (3.8, 3.8) {
    \textbf{\texttt{ContentRights-}}\\
    \textbf{\texttt{Precompile}}\\
    {\footnotesize EVM read-only (H160)}
};
%% Core pallet
\node[core, minimum width=5.5cm, minimum height=1.3cm] (pcr) at (0, 1.6) {
    \textbf{\texttt{pallet-content-rights}}\\[3pt]
    {\footnotesize 17 extrinsics \quad 10 storage maps \quad 22 events}
};
%% Underlying support pallets
\node[dep] (pnfts) at (-2.8, -0.8) {
    \textbf{\texttt{pallet-nfts}}\\
    {\footnotesize parent/child NFT storage}
};
\node[sib] (prv) at (2.8, -0.8) {
    \textbf{\texttt{pallet-rights-}}\\
    \textbf{\texttt{verifier}}\\
    {\footnotesize Merkle proof verification}
};
%% Arrows: callers to core
\draw[->, thick] (xcm.south) -- node[lbl, near start, left]{\texttt{xcm\_*} extrinsics} (pcr.north);
\draw[->, thick] (ink.south) -- node[lbl, near end, right]{forward calls} (pcr.north);
\draw[->, thick, dashed] (pre.south) -- node[lbl, near start, right]{storage reads} (pcr.north);
%% Arrows: core to dependencies
\draw[->, thick] (pcr.south) -- node[lbl, near end, left]{mint / nest / burn} (pnfts.north);
\draw[->, thick, dashed] (pcr.south) -- node[lbl, near end, right]{state for proofs} (prv.north);
\end{tikzpicture}
\caption{CCRMS implementation-level component diagram. Solid arrows indicate active extrinsic dispatch; dashed arrows indicate read-only storage access. \texttt{pallet-content-rights} is the sole authoritative state layer; the ink!\ contract and EVM precompile are stateless wrappers that forward calls or read storage, respectively. XCM delivers cross-chain calls via HRMP. \texttt{pallet-rights-verifier} enables remote parachains to validate rights state via Merkle storage proofs, without oracles or relayers.}
\label{fig:component-overview}
\end{figure}

The core of CCRMS is \texttt{pallet-content-rights}, a FRAME pallet that implements all business logic: content registration, the three monetization models (subscription, pay-per-view, and permanent ownership), on-chain royalty distribution to up to ten collaborators, automated subscription renewal via an \texttt{on\_initialize} hook, and cross-chain variants of every financial operation. It exposes 17 dispatchable extrinsics, 10 storage items, and 22 event variants; the full reference tables are in Appendix~\ref{app:code}. Because \texttt{check\_access} is invoked via XCM from remote parachains and each call is recorded on-chain for auditability, it is implemented as a signed extrinsic rather than a free runtime API. Content rights are represented as NFT collections using a parent--child nesting model built on \texttt{pallet-nfts}: a parent Content NFT is the canonical on-chain identity for each piece of content, and each user's access rights are a child NFT tagged with a \texttt{rights\_type} attribute (Subscription, PPV, or Ownership). The smart contract and precompile layers are stateless wrappers; all authoritative state resides in the pallet (Appendices~\ref{sec:contract-first-pivot} and~\ref{sec:ink-discontinuation}).

\section{Cross-Chain Integration}\label{sec:cross-chain-integration}

The core challenge in cross-chain integration is enabling remote parachains to acquire and verify content rights without a shared trust anchor or custodian. Two sub-problems follow: (1) how to dispatch authenticated, fee-paying rights operations across HRMP in a single atomic message; and (2) how to extend that reach to Ethereum without custodial bridges. The key lesson from both is that the CCRMS runtime's XCM configuration and payer--beneficiary decoupling are sufficient; no custom XCM extensions were required.

\subsection{XCM Message Flows}\label{subsec:xcm-flows}

Cross-chain operations use Polkadot XCM v3 to enable remote parachains to acquire content rights. The pallet exposes five XCM extrinsics: \texttt{xcm\_subscribe}, \texttt{xcm\_renew\_subscription}, \texttt{xcm\_purchase\_views}, \texttt{xcm\_purchase\_ownership}, and \texttt{xcm\_transfer\_ownership} (callable locally only after Finding~B remediation; see Appendix~\ref{app:security}). Each separates the dispatch origin (\textbf{payer}) from the \textbf{beneficiary} of the rights record: when a remote parachain sends a \texttt{Transact} message, the origin is the remote parachain's sovereign account, but the rights record is issued to a distinct \texttt{beneficiary} account representing the end user on the remote chain.

The runtime's XCM configuration follows the standard Cumulus pattern, with one customization: the \texttt{LocationToAccountId} chain adds \texttt{HashedDescription<\_,}\allowbreak \texttt{DescribeFamily<}\allowbreak\texttt{DescribeAllTerminal>{}>} and \texttt{GlobalConsensusParachainConvertsFor<\_,\,\_>} so the runtime can interpret Ethereum locations of the form \texttt{GlobalConsensus(Ethereum)} $\to$ \texttt{AccountKey20(0x...)}. This is required for Snowbridge (Section~\ref{subsec:snowbridge}); the full Ethereum-to-CCRMS flow is configured but has not been end-to-end tested.

Each XCM rights operation uses the three-instruction sequence \texttt{WithdrawAsset}, \texttt{BuyExecution}, and \texttt{Transact}. This is sufficient because rights operations are stateful side effects, not asset transfers (Appendix~\ref{app:code} shows the full message structure). The endianness and fee-sizing findings (A.5, A.6 in Table~\ref{tab:challenges}) are the two most operationally significant surprises in the XCM integration.

We measured end-to-end XCM latency at approximately 100 seconds of wall-clock time (a 3--5-block delta on the CCRMS chain) for \texttt{xcm\_subscribe}, \texttt{xcm\_purchase\_views}, and \texttt{xcm\_purchase\_ownership} on the local Zombienet topology, full measurements are reported in Chapter~\ref{chap:results}. The block delta underestimates the true elapsed time: the relay-chain hop and HRMP processing pipeline add $\sim$70 seconds of wall-clock overhead that is not reflected in the CCRMS block count. Polkadot~2.0 features (Async Backing, Agile Coretime, Elastic Scaling) would substantially reduce the latency floor in production.

A secondary pallet, \texttt{pallet-rights-verifier}, provides trustless verification of CCRMS rights status for other parachains via Merkle storage proofs. A remote parachain can verify an active subscription without oracles or relayers by submitting a proof against a known CCRMS state root, which the verifier replays against the Substrate trie. Three extrinsics (\texttt{verify\_ownership}, \texttt{verify\_subscription}, \texttt{verify\_view\_pack}) construct a storage key, validate the proof, and emit a result event. Proof validation costs roughly 100 million \texttt{ref\_time} per call and avoids any network round-trip.

\subsection{Snowbridge Integration with Ethereum}\label{subsec:snowbridge}

We extended the cross-chain architecture to Ethereum via Snowbridge, a trust-minimized bridge that uses light-client verification on both sides, eliminating custodians and relayer trust. During Phase~4 we bridged two 1-ETH transactions from a local Ethereum chain to a local AssetHub parachain using the Snowbridge pipeline, demonstrating that the CCRMS runtime can accept bridged Ethereum messages with the same XCM configuration as sibling-parachain messages (via \texttt{HashedDescription} and \texttt{GlobalConsensusParachainConvertsFor} in \texttt{LocationToAccountId}).

The topology spans four parachains and a local Ethereum network: CCRMS (parachain 100); Bridge Hub (parachain 1013), which hosts the inbound queue and the Ethereum beacon light client; AssetHub (parachain 1000), which holds the foreign-asset reserve; and the Rococo relay chain. Four HRMP channels link Bridge Hub $\leftrightarrow$ AssetHub and AssetHub $\leftrightarrow$ CCRMS, carrying bridge messages and enabling future Polkadot--Ethereum interoperability.

The Ethereum side runs Geth v1.17.1 (chain ID 11155111) as the execution layer, Lodestar v1.35.0 as the beacon node, and the 16-contract Snowbridge Gateway suite was deployed via Forge (a Solidity contract deployment tool from the Foundry framework). After initializing the beacon light client, two relayer processes run: a beacon relay that forwards Lodestar finality events to the \texttt{EthereumBeaconClient} pallet, and an execution relay that forwards Gateway events from Geth to the \texttt{EthereumInboundQueue} pallet.

The integration revealed three persistent issues, documented in Appendix~\ref{app:challenges}: a Lodestar version-pinning requirement that forces source compilation; a configuration bug in the fork version that produced invalid Merkle proofs; and a beacon state cache timing issue that required reordering the setup sequence.

The final step, using bridged Ether on AssetHub to pay for a CCRMS rights operation, is conceptually straightforward but was not implemented in the prototype. The CCRMS XCM side is already prepared; the remaining task is to configure \texttt{pallet-content-rights::PaymentCurrency} to accept both foreign Ether and the native token. This is listed as future work in Chapter~\ref{chap:conclusion}.

\section{Challenges and Solutions}\label{sec:challenges-solutions}

The CCRMS implementation encountered nine technical issues that shaped the final architecture and security posture. Table~\ref{tab:challenges} summarizes each issue, our response, and the key lesson; Appendix~\ref{app:challenges} provides full details.

\begin{longtable}{@{}L{0.05\linewidth} L{0.30\linewidth} L{0.30\linewidth} L{0.30\linewidth}@{}}
\caption{Implementation Challenges Summary}\label{tab:challenges}\\
\toprule
\textbf{\#} & \textbf{Challenge} & \textbf{Response} & \textbf{Key Lesson} \\
\midrule
\endfirsthead

\multicolumn{4}{l}{\emph{Table~\ref{tab:challenges} continued from previous page}}\\
\toprule
\textbf{\#} & \textbf{Challenge} & \textbf{Response} & \textbf{Key Lesson} \\
\midrule
\endhead

\midrule
\multicolumn{4}{r}{\emph{Continued on next page}}\\
\endfoot

\bottomrule
\endlastfoot

A.1 & RMRK 2.0 pallets frozen on SDK 0.9.36, incompatible with current monorepo & Replicated RMRK nesting semantics on \texttt{pallet-nfts} directly & Standards reimplementation is preferable when reference implementation is unmaintained \\
\addlinespace
A.2 & Contract-first design blocked: chain extensions non-functional for writes, debugging complexity & Inverted to pallet-first; retained ink! contract as thin wrapper & Pallet-first proved resilient to ink! discontinuation (A.9) \\
\addlinespace
A.3 & Snowbridge beacon relay generated proofs at wrong gindex (54 vs.\ 86) & \texttt{forkVersions} field means activation epoch, not version hex; changed to \texttt{\{"deneb": 0, "electra": 0\}} & Cryptographic systems fail silently on configuration errors \\
\addlinespace
A.4 & Lodestar v1.41.0 ignored \texttt{LODESTAR\_PRESET=mainnet} in dev mode (32 vs.\ 512 validators) & Pinned Lodestar v1.35.0, built from source & Light clients impose strict counterpart requirements; bridge toolchain compatibility is fragile \\
\addlinespace
A.5 & XCM tests: funds debited on source chain but \texttt{InsufficientPayment} on CCRMS & polkadot.js \texttt{toHex()} used big-endian; CCRMS expects little-endian; fixed with \texttt{toHex(true)} & Endianness bugs produce valid-looking but wrong addresses \\
\addlinespace
A.6 & XCM \texttt{BuyExecution} failed at 1B tokens; required ${\sim}$75--100B & \texttt{proof\_size} dominates \texttt{BlockRatioFee}, not \texttt{ref\_time} & Polkadot fee modeling is driven by proof size, not computation \\
\addlinespace
A.7 & \texttt{xcm\_transfer\_ownership} had no check that caller == \texttt{from} (Medium severity) & Added \texttt{ensure!(authorizer == from)}; unit test confirms block & Cross-chain extrinsics accepting identity parameters need explicit authorization \\
\addlinespace
A.8 & \texttt{MaxChildrenReached} at 50 subscribers per content item & Documented as known limitation; benchmark scripts rotate content items & Bounded FRAME collections trade benchmarkability for capacity limits \\
\addlinespace
A.9 & ink! development discontinued Jan 2026 & 95\% of logic in pallet; contract layer is non-load-bearing; migration deferred & Minimize dependency on non-essential components for ecosystem resilience \\
\end{longtable}

\section{Summary}\label{sec:chapter-5-summary}

Cross-chain rights operations are delivered via XCM v3 with payer--beneficiary decoupling, and the Snowbridge integration configured the path for Ethereum bridging, though the Ethereum-to-CCRMS flow has not been end-to-end tested. Nine challenges shaped the final architecture; the two most consequential---the RMRK-to-\texttt{pallet-nfts} pivot and the contract-first-to-pallet-first pivot---are examined in greater detail in Chapter~\ref{chap:discussion}. Full implementation details are provided in Appendix~\ref{app:code}.
